\documentclass{article}
\usepackage{algorithm, algorithmic, amsfonts, amsmath}

\begin{document}
\begin{algorithm}[!h]
\renewcommand{\algorithmicrequire}{\textbf{Input:}} % show \REQUIRE as Input
\renewcommand{\algorithmicensure}{\textbf{Output:}} % show \ENSURE as Output
	\caption{LUHC (Laplacian Unit-Hyperplane Classifier)} % algorithm title
	\begin{algorithmic}[1] % [1] numbers every line
	    \REQUIRE
	    $P;U;\nu\in(0,1]=$ upper bound on the fraction of positive margin errors; trade-off parameters $c_1,c_2$; graph parameters ($k$ nearest neighbors, bandwidth $\sigma$)
	    \ENSURE
	    Classifier $f$
	    \STATE $X\leftarrow P\cup U$
	    \STATE Build a $k$-NN similarity graph on $X$ with weight matrix $W$ ($W_{ij}=\exp(-\|x_i-x_j\|^2/\sigma^2)$ for neighbors, else 0)
	    \STATE $D\leftarrow \operatorname{diag}(\textstyle\sum_j W_{ij})$;\ $L\leftarrow D-W$ % graph Laplacian
	    \STATE Formulate the QPP for a decision unit-hyperplane $f(x)=w^{T}x+b$ with $\|w\|=1$: % a unit hyperplane, so f(x) is the signed distance
	    \STATE \quad maximize the margin $\rho$ by which positive examples lie on the positive side, with slack variables $\xi_i\geq 0$ on $P$ whose total fraction is bounded by $\nu$
	    \STATE \quad regularize with the manifold term $c_1\, \boldsymbol{f}^{T}L\boldsymbol{f}$, where $\boldsymbol{f}=(f(x_1),\dots,f(x_{|X|}))^{T}$ % geometric information: the decision values vary smoothly over the graph
	    \STATE \quad penalize with $c_2\sum_{x_j\in U} f(x_j)$ to push the unlabeled mass to the negative side % discriminative information: the unlabeled mass is pushed negative
	    \STATE Solve the dual QPP, whose size depends only on $|P|$ % a small QPP, so training is fast
	    \STATE $f(x)\leftarrow \operatorname{sign}(w^{T}x+b)$
		\RETURN $f$
	\end{algorithmic}
\end{algorithm}
% Shao et al., Information Sciences 314:152-168, 2015.
% The objective and its dual derivation are given in the original; this pseudocode records the structure only.
\end{document}
